

% NEEDS TO BE FILLED BY YOU

% Set your own title:
\newcommand{\titleoftheabstract}{The title of this abstract}

% Once your abstract is accepted, substitute [author name(s), affiliation(s) and e-mail(s)] with your author, affiliation and e-mail data:
\newcommand{\authornames}{Author name(s)}% for the pdf metadata
\newcommand{\authorsandaffiliations}{% for being displayed in the pdf
[author name(s), affiliation(s) and e-mail(s)]
% Example of a way to structure the author, affiliation and e-mail line:
% (Comment the line above and uncomment the line below to see how it looks like)
%Joachim Kokkelmans\hfill Free University of Bozen-Bolzano\hfill \href{mailto:jkokkelmans@unibz.it}{jkokkelmans@unibz.it}%
}

% END OF: NEEDS TO BE FILLED BY YOU


% MfM guidelines set 12pt as a minimum font size:
\documentclass[12pt]{article}

% Choosing a font which supports IPA:
% Doulos SIL supports IPA but not bold or italic; an alternative for bold and italic text is provided with Times New Roman, which does not support IPA however
% IPA- and bold/italic-compatible fonts which can be used instead:
% Charis SIL, CMU Bright, CMU Concrete, CMU Sans Serif, CMU Serif, Cousine, DejaVu Sans, Tinos (almost full IPA), Noto Sans (partially), Noto Serif (partially), Arimo (partially), Cardo (partially), DejaVu Serif (partially), EB Garamond (partially), Fira Sans (partially), Garamond Libre (partially), Inter (partially), Junicode (partially), Lato (partially), Libertinus Sans (partially), Libertinus Serif (partially), STIX (partially)
\usepackage{fontspec}
\setmainfont[
 BoldFont={Times New Roman Bold},
 ItalicFont={Times New Roman Italic},
 BoldItalicFont={Times New Roman Bold Italic}
 ]{Doulos SIL}

% For bibliographic references:
% Uses the Linguistic Society of America (LSA)-based sp citation style from S&P. Last names such as 'van Beethoven' will be ranked alphabetically as B-initial unless they are put in between curly brackets in references.bib
\usepackage[sectionbib]{natbib}
\bibpunct{(}{)}{;}{a}{}{,}
\setcitestyle{notesep={: }}
\usepackage{bibentry}% allows for the use of \nobibliography{}

% For the margins on the sides of the page:
% 2.5cm = 0.984252 inch, so 2.5cm allows for smaller margins than 1 inch
% A4 allows for a larger page area than US letter, so A4 is chosen
\usepackage[a4paper, margin=2.5cm]{geometry}

% Allows to include images and to wrap them next to the flowing text:
% Images must be situated in the folder Images
\usepackage{graphicx, wrapfig}
\graphicspath{ {./Images/} }

% Allows to use colours:
% Sets the colour of links within the document and towards web pages
\usepackage{xcolor}
\definecolor{customcol}{RGB}{33, 45, 141}

% For linguistic glosses and examples:
\usepackage{gb4e}

% To put text in two columns:
\usepackage{multicol}

% This defines the environment itemize* in which the spacing between items can be customised (default = 1pt):
\newenvironment{itemize*}%
  {\begin{itemize}%
    \setlength{\itemsep}{2pt}%
    \setlength{\parskip}{1pt}}%
  {\end{itemize}}

% Allows to have interactive links within the document and towards web pages:
\usepackage{url}
\usepackage{hyperref}
\hypersetup{linktoc=all, hidelinks, colorlinks=true, linkcolor=black, anchorcolor=customcol, citecolor=black, filecolor=customcol, menucolor=customcol, runcolor=customcol, urlcolor=customcol, pdftitle={\titleoftheabstract{}}, pdfauthor={\authornames{}}}


\begin{document}

\thispagestyle{empty}
\begin{center}
    {\bfseries \titleoftheabstract{}}
    \\\vspace{0.16cm}
    \authorsandaffiliations{}
\end{center}\vspace{-0.16cm}

% -----------------------------------------------------------------------
% -----------------------------------------------------------------------
% -----------------------------------------------------------------------

% Here the actual content of your abstract starts:
% Leave one or more empty row(s) between paragraphs to start a new paragraph with left indentation; if you wish to avoid the left indentation, add \noindent to the beginning of the new paragraph.


% If you don't want clickable interactive links (e.g. the author name and year in the in-text bibliographical references), uncomment \begin{NoHyper} here and \end{NoHyper} at the end of the text below.
%\begin{NoHyper}

\noindent This template follows the guidelines for abstract submission to the \href{http://www.lel.ed.ac.uk/mfm/mfm.html}{Manchester Phonology Meeting (MfM)}:
\begin{multicols}{2}
    \begin{itemize*}
        \item 1 page maximum
        \item A4 paper
        \item 2.5cm margins
        \item single spacing
        \item minimum 12pt font size
        \item normal character spacing
        \item optional full bibliography
        \item IPA support
    \end{itemize*}
\end{multicols}

\noindent With the Doulos SIL or Charis SIL fonts, you can use IPA characters such as [ɑ æ ɐ ɑ̃ β ɓ ʙ ç ɕ ð d͡ʒ ɖ ɗ ə ɚ ɵ ɘ ɛ ɜ ɝ ɛ̃ ɞ ɠ ɢ ʛ ħ ɦ ɥ ɧ ʜ ɪ ɨ ɪ ̈ ʝ ɟ ʄ ɫ ɭ ɬ ʟ ɮ ɱ ŋ ɲ ɳ ɴ ɔ œ ɒ ɔ̃ ɶ ø ɸ ɾ ɹ ʁ ʀ ɻ ɽ ɺ ʃ ʂ θ t͡ʃ t͡s ʈ ʊ ʉ ʌ ʋ ⱱ ɯ ʍ ɰ χ ɣ ʎ ʏ ɤ ʒ ʐ ʑ ʔ ʕ ʡ ʢ | ‖ ∅ → ˈ ˌ ː ʰ ʲ ʷ ̚ ˩ ˨ ˧ ˦ ˥ ˥ ˩ ˩ ˥ ˩ ˨ ˦ ˥ ˧ ˦ ˧ ˩˨˧˦˥˥˩˩˥˩˨˦˥˧˦˧ ˤ ˀ ᵝ ᵊ ʱ ˡ ⁿ ʳ ᵗ ˠ ʼ ǀ ǁ ǃ ǂ ʘ ʩ ʪ ʫ ꞎ ʬ ʭ ¡ ʶ ⁼ ↓ ↀ].\\

\noindent Examples can be provided like this:
\begin{exe}
    \ex Types of anchoring
    \begin{xlista}
        \ex Left-oriented
        \ex Right-oriented
        \ex Stress-dependent
    \end{xlista}
\end{exe}

\noindent Glosses can be provided like this:% Cfr. https://www1.essex.ac.uk/linguistics/external/clmt/latex4ling/examples/gb4e/
\begin{exe}
    \ex Example of \emph{h-aspiré} in standard French:
    \begin{xlisti}
        \ex[*]{Un hacker [ɛ̃.na.kœʀ] sait utiliser un ordinateur.}
        \ex[]{Un hacker [ɛ̃.a.kœʀ] sait utiliser un ordinateur.}
        \exi{}{\gll Un hacker sait utiliser un ordinateur.\\
        A hacker can use a computer\\}
    \end{xlisti}
\end{exe}

Here is how to cite \citet[17]{Chomsky1965} with parentheses around the year (and optionally the page, if provided as argument), \citep[68]{Chomsky1965} with parentheses around everything, or \citealt{Chomsky1965} without parentheses. An author entry \{van Beethoven, Ludwig\} in \emph{references.bib} is alphabetically ranked in the bibliography as starting with a \textbf{B}, but the author %
%
\begin{wrapfigure}{r}{4cm}
    \centering
    \fbox{\includegraphics[width=4cm]{Images/28mfmlogo.jpeg}}
    \caption{The logo for the 28\textsuperscript{th} MfM edition.}
    \label{fig:MfMLogo}
\end{wrapfigure}
%
entry \{\{van Beethoven\}, Ludwig\} is ranked as starting with a \textbf{V} \citep{vanBeethoven1824}.\\

\noindent An image without a caption can be added with the simple command \emph{includegraphics}; an image with a caption can be added with the same command within a \emph{figure} environment; and an image placed left or right from the main text (as in this example) can be generated by using the same configuration but with a \emph{wrapfigure} environment (and specifying the alignment and width of the image).\\\\\\

%\end{NoHyper}


% Here the actual content of your abstract stops.
% -----------------------------------------------------------------------
% -----------------------------------------------------------------------
% -----------------------------------------------------------------------


\renewcommand{\bibsection}{}% Sets an empty bibliography title
\setlength{\bibsep}{0.0pt}% Sets normal paragraph spacing in the bibliography
\bibliographystyle{sp}
% If you wish to display the references at the bottom of the page, the next line must have '\bibliography', otherwise it must have '\nobibliography'
\bibliography{references}

\end{document}